% ======================================================================
% 6. přednáška (7. 11. 2022) – Kosmické záření (pokr.), biologické účinky,
% symetrie a spin, základní členění částic
% ======================================================================
\section{Kosmické záření, biologické účinky záření a symetrie}

\strana{P06-1}
\subsection{Kosmické záření – pokračování}

Minule jsme odhadli, že mion ztratí průchodem atmosférou asi \(2\,\text{GeV}\).
Mion je ale nestabilní (\(\mu^-\to e^-\bar\nu_e\nu_\mu\)) se střední dobou
života \(\tau \approx 2{,}2\,\mu\text{s}\). Klasicky by za tuto dobu urazil
nanejvýš \(c\tau\):
\[
  c\tau(\mu) = 658\,\text{m}, \qquad m_\mu = 106\,\text{MeV}.
\]
Miony vznikají ve výšce kolem \(15\,\text{km}\) – bez relativity by se tedy
k zemi nedostal téměř žádný. Že je přesto pozorujeme, je klasický důkaz
\kw{dilatace času}: v laboratorní soustavě žije mion \(\gamma\)-krát déle.

Střední doletová vzdálenost (\kw{střední dráha do rozpadu}) v laboratoři je
\[
  l = \gamma(\beta c)\tau = \gamma\beta\,(c\tau)
    = \frac{E}{m}\,\frac{p}{E}\,c\tau
    = \frac{\sqrt{E^2-m^2}}{m}\,c\tau = 37{,}4\,\text{km}.
\]
Využili jsme \(\gamma = E/m\) a \(\beta = p/E\) (jednotky \(c=1\)), takže
\(\gamma\beta = p/m\). Hodnota \(37{,}4\,\text{km}\) odpovídá
\(p/m = 37{,}4/0{,}658 \approx 57\), tj. \(p \approx 6\,\text{GeV}\) – to je
energie, kterou mion potřebuje, aby po ztrátě \(\sim 2\,\text{GeV}\) dorazil
na zem s typickou energií \(\sim 4\,\text{GeV}\).

\begin{priklad}[Jaká část mionů doletí?]
Rozpad je exponenciální v dráze stejně jako v čase: pravděpodobnost, že se
mion rozpadne na dráze \(L\), je \(1 - e^{-L/l}\). Pro \(L = 15\,\text{km}\):
\[
  1 - e^{-\frac{15\,\text{km}}{37{,}4\,\text{km}}} \doteq \frac13
  \quad\Longrightarrow\quad \frac23 \text{ mionů doletí k povrchu Země.}
\]
(\(e^{-0{,}40}\approx 0{,}67\).) Odhad zanedbává, že mion cestou zpomaluje
(\(\gamma\) klesá), takže skutečná ztráta rozpadem je o něco větší; řádově
je však výsledek správný.
\end{priklad}

\subsection{Biologické účinky záření}

Ionizující záření poškozuje tkáň tím, že v ní ukládá energii – ionizuje
molekuly (především vodu a DNA). Základní měřítko je proto deponovaná
energie na jednotku hmotnosti; různé druhy záření ale při stejné deponované
energii škodí různě, a proto se zavádí několik druhů dávek.

\begin{itemize}
  \item \kw{Dávka} (absorbovaná dávka)
  \[
    \vzorec{1\,\text{Gy}\ (\text{gray}) = 1\,\text{J}/\text{kg}}
  \]
  – čistě fyzikální veličina: energie pohlcená v kilogramu látky.
  \item \kw{Ekvivalentní dávka} zohledňuje druh záření pomocí radiačního
  váhového faktoru \(w_R\):
  \[
    \vzorec{1\,\text{Sv} = 1\,\text{Gy}\times\text{váha}}
  \]
  \begin{itemize}
    \item \(\gamma\), \(e^\pm\): váha \(=1\),
    \item \(p\), \(\pi^\pm\): váha \(=2\),
    \item \(\alpha\), štěpné fragmenty, těžké ionty: váha \(=20\),
    \item neutrony: váha závisí na energii, maximální váha \(=20\)
    (kolem \(1\,\text{MeV}\)).
  \end{itemize}
  Proč? Fotony a elektrony ionizují „řídce`` (malá lineární hustota
  ionizace), takže poškození DNA jsou izolovaná a buňka je umí opravit.
  Těžké, pomalé, mnohonásobně nabité částice (\(\alpha\), fragmenty) ukládají
  energii velmi hustě podél krátké stopy a způsobují vícečetné zlomy DNA,
  které jsou opravitelné jen obtížně. Neutrony ionizují nepřímo – přes
  odražené protony a jaderné reakce – proto jejich váha závisí na energii.
  \item \kw{Efektivní dávka} – součet ekvivalentních dávek jednotlivých
  orgánů vážený tkáňovými váhovými faktory \(w_T\) (různé orgány mají různou
  citlivost, např. kostní dřeň vs.\ kůže), \(E = \sum_T w_T H_T\). Udává se
  rovněž v Sv a vyjadřuje celkové riziko pro organismus při nerovnoměrném
  ozáření.
\end{itemize}

\kw{Účinky} záření dělíme na
\begin{itemize}
  \item \kw{stochastické} – náhodné, bez prahu; s dávkou roste
  \emph{pravděpodobnost} (např. vzniku nádoru nebo genetické změny), nikoli
  závažnost;
  \item \kw{deterministické} – od \(\sim 100\,\text{mGy}\); nastávají vždy po
  překročení prahu a jejich závažnost roste s dávkou:
  \begin{itemize}
    \item lokální poškození buněk (tkáně),
    \item dávka \(\gg\) přírodní pozadí (\(3\)–\(4\,\text{mSv}/\text{rok}\)).
  \end{itemize}
\end{itemize}

\strana{P06-2}
\kw{Pásma dávek} (pro představu):
\begin{itemize}
  \item vysoké: \(> 1\,\text{Gy}\) (hasiči v Černobylu),
  \item střední: \(100\,\text{mGy}\)–\(1\,\text{Gy}\) (sanační práce v Černobylu),
  \item nízké: \(10\)–\(100\,\text{mGy}\) (vícefázová CT, vyšetření CT
  \(50\)–\(70\,\text{mGy}\)) \pozn{čtení „50–70`` nejisté},
  \item velmi malé: \(< 10\,\text{mGy}\) (běžný rentgenový snímek).
\end{itemize}
Letadlo Paříž–New York: \(0{,}05\,\text{mSv}\) – ve výšce \(\sim 10\,\text{km}\)
nás atmosféra chrání před kosmickým zářením méně. Jeden let tedy odpovídá
zhruba několika dnům přírodního pozadí.

\subsubsection{Použití v medicíně}
\begin{itemize}
  \item \kw{PET} (pozitronová emisní tomografie) – pacientovi se podá
  \(\beta^+\) zářič (např. \({}^{18}\)F v glukóze); pozitron anihiluje
  s elektronem a dva fotony \(511\,\text{keV}\) letící opačnými směry
  určí přímku, na níž k anihilaci došlo (Davídkovy slajdy „Aplikace
  v medicíně``),
  \item \kw{rtg} – zobrazení díky různé absorpci (fotoefekt \(\sim Z^5\)).
\end{itemize}

\kw{Léčení} – ozařování nádorů. Klasickým zdrojem je kobalt:
\[
  {}^{60}\text{Co} \to {}^{60}\text{Ni}^* + e^- + \bar\nu_e \to {}^{60}\text{Ni} + 2\gamma .
\]
Rozpad \(\beta^-\) vede do excitovaného stavu niklu, který se kaskádou
dvou fotonů (\(1{,}17\) a \(1{,}33\,\text{MeV}\)) deexcituje do základního
stavu. Pronikavé fotony o energii \(\sim\)\,MeV jsou to, co se k ozařování
používá (\(T_{1/2}\approx 5{,}3\) roku).

\begin{center}
\begin{tikzpicture}[>=Stealth, thick]
  \draw (0,3) -- (1.6,3) node[midway, above] {\({}^{60}\)Co};
  \draw[->] (1.2,3) -- (3.0,2.2) node[midway, above right] {\(\beta^-\)};
  \draw (2.6,2.2) -- (4.6,2.2) node[right] {\(2{,}51\,\text{MeV}\)};
  \draw (2.6,1.0) -- (4.6,1.0) node[right] {\(1{,}33\,\text{MeV}\)};
  \draw (2.6,0) -- (4.6,0) node[right] {\(0\)\quad \({}^{60}\)Ni};
  \draw[decorate,decoration={snake,amplitude=2pt,segment length=6pt}] (3.6,2.2) -- (3.6,1.15);
  \draw[->] (3.6,1.15) -- (3.6,1.0);
  \node[left] at (3.5,1.6) {\(\gamma\)};
  \draw[decorate,decoration={snake,amplitude=2pt,segment length=6pt}] (3.6,1.0) -- (3.6,0.15);
  \draw[->] (3.6,0.15) -- (3.6,0);
  \node[left] at (3.5,0.5) {\(\gamma\)};
\end{tikzpicture}
\end{center}
(Energie hladin jsou doplněny z tabulek, v náčrtku v poznámkách nejsou.)

\kw{PTC} (Proton Therapy Center) – protony až do \(250\,\text{MeV}\).
Výhoda protonů: ionizační ztráty rostou se zpomalováním (\(\sim\beta^{-2}\)),
takže většinu energie uloží na konci dráhy v tzv. \kw{Braggově píku}.
Volbou energie lze pík umístit přesně do nádoru a okolní zdravou tkáň
šetřit; \(250\,\text{MeV}\) odpovídá doletu kolem \(35\,\text{cm}\) ve vodě
(tkáni).

\subsubsection{Dozimetry}
Dozimetry měří dávku, kterou obdržel jejich nositel:
\begin{itemize}
  \item \kw{filmové} – zčernání fotografické emulze,
  \item \kw{termoluminiscenční} – LiF, CaF\(_2\), MgBeO\(_4\)
  \pozn{poslední vzorec nejistý; obvyklé TLD materiály jsou např. BeO nebo
  MgB\(_4\)O\(_7\)}: záření převede elektrony do metastabilních pastí,
  při zahřátí se uvolní a vyzáří světlo úměrné dávce,
  \item \kw{elektronické} – Si (polovodičová dioda) nebo Geigerův–Müllerův
  počítač; umožňují okamžitý odečet.
\end{itemize}

\begin{priklad}[Dávka od drobného zdroje \({}^{137}\)Cs]
Uvažujme \oprava{\({}^{137}\)Cs}{\({}^{134}\)Cs} o aktivitě \(363\,\text{kBq}\) – drobný (kalibrační)
zdroj ionizujícího záření. Rozpadá se \(\beta^-\) na baryum: v \(95\,\%\)
případů do excitovaného stavu, který vyzáří foton
\(E_\gamma = 662\,\text{keV}\), v \(5\,\%\) přímo do základního stavu.

\begin{center}
\begin{tikzpicture}[>=Stealth, thick]
  \draw (0,2.6) -- (2.2,2.6) node[midway, above] {\({}^{137}\)Cs};
  \draw[->] (1.5,2.6) -- (3.4,1.45) node[midway, above right] {\(\beta^-\ (95\,\%)\)};
  \draw[->] (1.5,2.6) -- (3.4,0.05) node[pos=0.6, left] {\(\beta^-\ (5\,\%)\)};
  \draw (3.0,1.4) -- (4.8,1.4) node[right] {Ba};
  \draw (3.0,0) -- (4.8,0);
  \draw[decorate,decoration={snake,amplitude=2pt,segment length=6pt}] (4.0,1.4) -- (4.0,0.15);
  \draw[->] (4.0,0.15) -- (4.0,0);
  \node[right] at (4.1,0.7) {\(\gamma\), \(E_\gamma = 662\,\text{keV}\)};
\end{tikzpicture}
\end{center}

Jakou dávku za rok obdrží člověk (\(80\,\text{kg}\)), který je ve vzdálenosti
\(0{,}5\,\text{m}\) od zdroje? Předpokládáme: zdroj září izotropně, člověk
nastavuje plochu \(\sim 1\,\text{m}^2\), takže zachytí podíl
\(1\,\text{m}^2/(4\pi\,(0{,}5\,\text{m})^2)\) všech fotonů, a každý
zachycený foton v těle uloží celou energii \(662\,\text{keV}\) (elektrony
z \(\beta\) rozpadu se zastaví už v obalu zdroje/vzduchu). Rok je
\(\approx \pi\times 10^7\,\text{s}\).
\begin{align*}
  D &= \frac{1\,\text{m}^2}{4\pi\, 0{,}5^2\,\text{m}^2}\cdot 363\,\text{kBq}
       \cdot 662\,\text{keV} \,/\, 80\,\text{kg} = \\
    &= \frac{1}{\pi}\; 363\times10^{3}\,\text{s}^{-1}\;
       \pi\, \oprava{10^{7}}{10^{4}}\,\text{s}/\text{rok}
       \cdot 662\times10^{3}\,\text{eV}\times 1{,}6\times10^{-19}\,\text{J}/\text{eV}
       \,/\, 80\,\text{kg} = \\
    &= 4{,}8\times10^{-3}\,\text{Gy}/\text{rok}.
\end{align*}
Výsledek \(\approx 5\,\text{mGy}/\text{rok}\) (pro fotony \(\approx 5\,\text{mSv}/\text{rok}\))
je srovnatelný s přírodním pozadím \(3\)–\(4\,\text{mSv}/\text{rok}\) – takový
zdroj tedy nelze nosit trvale v kapse, ale krátkodobá manipulace je
bezpečná. (Výsledek \(4{,}8\times10^{-3}\) vychází jen s
\(\pi\times10^{7}\,\text{s}\); s \(10^4\) by byl o tři řády menší.)

Zbývá ověřit, že fotony nepohltí vzduch mezi zdrojem a člověkem.
Radiační délka vzduchu (PDG):
\[
  X_0(\text{air}) = 36{,}6\,\text{g}\,\text{cm}^{-2}, \qquad
  \rho(\text{air}) = 1{,}2\,\text{kg}\cdot\text{m}^{-3}
  \quad\Rightarrow\quad X_0 = \frac{36{,}6\,\text{g}\,\text{cm}^{-2}}{1{,}2\times10^{-3}\,\text{g}\,\text{cm}^{-3}}
  \approx 300\,\text{m},
\]
\[
  e^{-\frac{0{,}5}{300}} \approx 1 .
\]
Zeslabení na půl metru vzduchu je tedy zanedbatelné.
\end{priklad}

\begin{vysvetleni}[Radiační délka vs.\ absorpční délka]
\(X_0\) je přesně řečeno charakteristická délka pro brzdné záření elektronů
a tvorbu párů (střední volná dráha fotonu vůči tvorbě párů je
\(\tfrac97 X_0\)). Pro fotony \(662\,\text{keV}\) dominuje Comptonův rozptyl
a správnější je použít hmotnostní součinitel zeslabení
\(\mu/\rho \approx 0{,}077\,\text{cm}^2/\text{g}\), což dává
\(\lambda \approx 1/(0{,}077\cdot 1{,}2\times10^{-3})\,\text{cm} \approx 110\,\text{m}\).
Závěr se nemění: \(e^{-0{,}5/110} \approx 0{,}995 \approx 1\).
\end{vysvetleni}

\strana{P06-3}
\subsection{Pohádka o symetriích}

Proč se v jaderné a částicové fyzice tolik mluví o symetriích? Podle
\kw{Noetherové věty} odpovídá každé spojité symetrii zákon zachování; a
kvantová čísla, kterými částice klasifikujeme (spin, izospin, parita, \dots),
jsou právě „štítky`` reprezentací těchto symetrií. Než přejdeme ke
klasifikaci částic, ukážeme si, jak symetrie vede na zachovávající se
veličiny a jak z chování vlnové funkce při rotacích vyplývá spin.

\begin{itemize}
  \item \kw{Izospin} \(p\)–\(n\): silná interakce „nevidí`` rozdíl mezi
  protonem a neutronem, proto je chápeme jako dva stavy jedné částice
  (nukleonu) s izospinem \(\tfrac12\) – podobně jako dva stavy spinu.
  Deuteron je izosinglet:
  \[
    I({}^2_1\text{H}) = 0 .
  \]
  \item \kw{Zachovávající se veličina} \(\hat z\): pozorovatelná se zachovává,
  pokud komutuje s hamiltoniánem,
  \[
    \vzorec{[\hat H, \hat z] = 0} .
  \]
  Z Heisenbergovy pohybové rovnice
  \(\frac{\dd}{\dd t}\langle \hat z\rangle = \frac{i}{\hbar}\langle[\hat H,\hat z]\rangle\)
  plyne, že její střední hodnota se v čase nemění. Je-li \(\hat z\)
  generátorem nějaké transformace (viz níže), znamená tato komutace přesně
  to, že hamiltonián je vůči transformaci symetrický.
\end{itemize}

\subsubsection{Translace}
Posuňme vlnovou funkci o \(\varepsilon\) a rozviňme do Taylorovy řady:
\begin{align*}
  \psi(x+\varepsilon) &= \psi(x) + \varepsilon\,\pdv{\psi}{x}
     + \frac{1}{2!}\,\varepsilon^2\,\pdv[2]{\psi}{x} + \dots
   = e^{\varepsilon\pdv{x}}\,\psi(x)
   = e^{\frac{i\varepsilon\left(-i\hbar\pdv{x}\right)}{\hbar}}\,\psi(x) = \\
  &= e^{i\varepsilon\hat p_x/\hbar}\,\psi(x),
\end{align*}
a ve třech rozměrech
\[
  \vzorec{\psi(\vec x + \vec\varepsilon) = e^{i\vec\varepsilon\cdot\hat{\vec p}/\hbar}\,\psi(\vec x)} .
\]
Taylorova řada je formálně exponenciála operátoru derivace; ten po
vynásobení \(-i\hbar\) je právě operátor hybnosti. \kw{Hybnost je generátor
translací}. Je-li hamiltonián invariantní vůči posunutí (homogenní prostor),
komutuje s \(\hat{\vec p}\) a hybnost se zachovává.

\subsubsection{Rotace}
Rotace o úhel \(\alpha\) kolem osy \(z\) působí na souřadnice maticí
\[
  R_{z,\alpha} = \begin{pmatrix} \cos\alpha & \sin\alpha & 0\\ -\sin\alpha & \cos\alpha & 0\\ 0 & 0 & 1\end{pmatrix}.
\]
Stejně jako u translací nás zajímají \emph{infinitezimální} rotace
(\(\cos\delta\alpha\approx 1\), \(\sin\delta\alpha\approx\delta\alpha\)) –
z nich lze konečné rotace složit. Rozepíšeme je jako jednotkovou matici
plus malá oprava, jejíž koeficient definuje generátor \(G_k\):
\[
  R_{z,\delta\alpha} = \mathbf{1} + \delta\alpha\,
  \underbrace{\begin{pmatrix} 0 & 1 & 0\\ -1 & 0 & 0\\ 0 & 0 & 0\end{pmatrix}}_{iG_3},
  \qquad
  R_{y,\delta\alpha} = \mathbf{1} + \delta\alpha\,
  \underbrace{\begin{pmatrix} 0 & 0 & -1\\ 0 & 0 & 0\\ 1 & 0 & 0\end{pmatrix}}_{iG_2},
\]
\[
  R_{x,\delta\alpha} = \mathbf{1} + \delta\alpha\,
  \underbrace{\begin{pmatrix} 0 & 0 & 0\\ 0 & 0 & 1\\ 0 & -1 & 0\end{pmatrix}}_{iG_1}.
\]
\strana{P06-4}
Rotace o konečný úhel \(\alpha\) kolem obecné osy s jednotkovým vektorem
\(\vec n\) vznikne opakováním mnoha infinitezimálních
(\(\lim_{N\to\infty}(\mathbf{1} + i\vec G\cdot\vec n\,\alpha/N)^N\)):
\[
  \vzorec{R_{\vec n,\alpha} = e^{i(\vec G\cdot\vec n)\alpha}} .
\]
Generátory nekomutují (rotace kolem různých os nejsou záměnné), jejich
komutační relace
\[
  \vzorec{[\hat G_i, \hat G_j] = i\,\varepsilon_{ijk}\,\hat G_k}
\]
(\(\varepsilon_{ijk}\) je Levi-Civitův symbol) definují \kw{Lieovu algebru}
grupy rotací. Jsou to přesně komutační relace momentu hybnosti
(dělené \(\hbar\)), což naznačuje, že \kw{moment hybnosti je generátor
rotací}. Ověříme to přímo.

Operátor orbitálního momentu hybnosti je
\[
  \hat L_i = \varepsilon_{ijk}\,\hat x_j\,\hat p_k
  \quad\Rightarrow\quad
  \underline{L_3} = \hat x\hat p_2 - \hat y\hat p_1
  = -i\hbar\left[x\pdv{y} - y\pdv{x}\right]
  = \underline{-i\hbar\pdv{\varphi}} .
\]
Poslední rovnost ověříme pomocí sférických souřadnic (v konvenci, kde
\(\theta\) je zeměpisná šířka):
\[
  x = r\cos\theta\cos\varphi, \qquad
  y = r\cos\theta\sin\varphi, \qquad
  z = r\sin\theta .
\]
Řetízkovým pravidlem
\begin{align*}
  \pdv{\varphi} &= \pdv{x}\pdv{x}{\varphi} + \pdv{y}\pdv{y}{\varphi}
                 + \pdv{z}\pdv{z}{\varphi} = \\
  &= \pdv{x}\,r\cos\theta(-\sin\varphi) + \pdv{y}\,r\cos\theta\cos\varphi + 0 = \\
  &= -y\pdv{x} + x\pdv{y} .
\end{align*}
Nyní můžeme zopakovat trik s Taylorovou řadou pro úhel \(\varphi\):
\[
  \psi(\varphi+\varepsilon) = e^{\varepsilon\pdv{\varphi}}\,\psi(\varphi)
  = e^{i\varepsilon\hat L_3/\hbar}\,\psi(\varphi),
\]
neboť \(\pdv{\varphi} = i\hat L_3/\hbar\). Pro obecnou rotaci kolem osy
\(\vec n\):
\[
  \to\ \text{pro obecnou rotaci:}\quad e^{i(\hat{\vec L}\cdot\vec n)\varepsilon/\hbar} .
\]
Srovnáním s \(R_{\vec n,\alpha} = e^{i(\vec G\cdot\vec n)\alpha}\) vidíme, že
pro skalární vlnovou funkci \(\vec G = \hat{\vec L}/\hbar\).

\subsubsection{Se spinem}
Pokud vlnová funkce má více složek, rotace míchá i tyto složky mezi sebou
– a generátorem je pak celkový moment hybnosti:
\[
  \vec J = \vec L + \vec S : \qquad \vec G \to \vec G = (\vec L + \vec S)/\hbar .
\]
Spin \(\vec S\) je tedy ta část generátoru rotací, která působí na
„vnitřní`` (složkové) stupně volnosti. Podle toho, jak se vlnová funkce
transformuje, rozlišujeme:

\paragraph{Skalární funkce.}
\[
  \hat R_{\vec n,\alpha}\,\psi(\vec x) = \psi\big(\underbrace{\hat R^{-1}_{\vec n,\alpha}\,\vec x}_{\vec x'}\big),
  \qquad \vec J = \vec L + \vec S,\ \ \vec S = 0 .
\]
Rotuje se pouze argument (hodnota funkce v novém bodě je hodnota ve
„zpětně otočeném`` bodě), funkce sama je jediné číslo. Popisuje
\pozor{částice se spinem 0} (např. piony).

\paragraph{Vektorová funkce.}
\[
  \psi(\vec x) = \begin{pmatrix} \psi_1(\vec x)\\ \psi_2(\vec x)\\ \psi_3(\vec x)\end{pmatrix},
  \qquad
  \hat R_{\vec n,\alpha}\,\psi(\vec x)
  = \underbrace{\hat R_{\vec n,\alpha}\begin{pmatrix}\psi_1\\ \psi_2\\ \psi_3\end{pmatrix}}_{\psi'}
    \big(\underbrace{\hat R^{-1}_{\vec n,\alpha}\,\vec x}_{\vec x'}\big)
  = \oprava{e^{i(\hat{\vec S}\cdot\vec n)\alpha/\hbar}}{e^{i(\hat{\vec S}\cdot\vec n)/\hbar}}\,\psi\big(\hat R^{-1}\vec x\big).
\]
Teď se otáčí i samotný vektor \((\psi_1,\psi_2,\psi_3)\) – stejnou maticí jako
souřadnice. Generátory této vnitřní rotace jsou (\(\hat S_k = \hbar\hat G_k\)):
\[
  \hat S_1 = \hbar\hat G_1 = \hbar\begin{pmatrix} 0&0&0\\ 0&0&-i\\ 0&i&0\end{pmatrix},\quad
  \hat S_2 = \hbar\begin{pmatrix} 0&0&+i\\ 0&0&0\\ -i&0&0\end{pmatrix},\quad
  \hat S_3 = \hbar\begin{pmatrix} 0&-i&0\\ i&0&0\\ 0&0&0\end{pmatrix}.
\]
Snadno ověříme (např. \(\hat S_3^2 = \hbar^2\,\mathrm{diag}(1,1,0)\) a cyklicky):
\[
  \hat S_1^2 + \hat S_2^2 + \hat S_3^2 = 2\hbar^2\,\mathbf{1}
  = \oprava{s(s+1)\hbar^2}{s(s+1)}\,\mathbf{1}
  \quad\Longrightarrow\quad \text{spin} = 1 .
\]
Vektorové pole (např. vektorový potenciál fotonu \(\vec A\)) tedy popisuje
\pozor{částice se spinem 1}.

\strana{P06-5}
\paragraph{Spinorová funkce.}
\[
  \psi(\vec x) = \begin{pmatrix} \psi_\uparrow(\vec x)\\ \psi_\downarrow(\vec x)\end{pmatrix},
  \qquad
  \hat R_{\vec n,\alpha}\,\psi(\vec x)
  = \underbrace{e^{i(\hat{\vec S}\cdot\vec n)\alpha/\hbar}\,\psi}_{\psi'}
    \big(\underbrace{\hat R^{-1}_{\vec n,\alpha}\,\vec x}_{\vec x'}\big),
\]
kde pro dvousložkovou funkci \(\hat{\vec S} = \frac{\hbar}{2}\vec\sigma\)
(\(\vec\sigma\) jsou Pauliho matice). Protože \((\vec\sigma\cdot\vec n)^2 = \mathbf{1}\),
lze exponenciálu sečíst stejně jako \(e^{ix} = \cos x + i\sin x\):
\[
  \vzorec{\hat S_{\vec n,\alpha} = \Big(\cos\frac{\alpha}{2}\Big)\mathbf{1}
  + i\sin\frac{\alpha}{2}\,(\vec\sigma\cdot\vec n)}
  \qquad\Longrightarrow\qquad \text{spin } \tfrac12 .
\]
Charakteristické jsou poloviční úhly: při otočení o \(\alpha = 2\pi\) dostaneme
\(\hat S = -\mathbf{1}\), spinor změní znaménko; teprve otočení o \(4\pi\) jej
vrátí do původního stavu. Takto se transformují elektrony, nukleony,
kvarky – \pozor{částice se spinem \(\tfrac12\)}.

\begin{vysvetleni}[Shrnutí: spin jako vlastnost transformace]
Spin tedy není „rotace kuličky kolem osy``, ale údaj o tom, \emph{kolik
složek} má vlnová funkce a jakou maticí se tyto složky při otočení míchají:
1 složka \(\to\) spin 0, 3 složky (vektor) \(\to\) spin 1, 2 složky (spinor)
\(\to\) spin \(\tfrac12\). Obecně má reprezentace se spinem \(s\) celkem
\(2s+1\) složek. Tento pohled vysvětluje, proč je spin čistě kvantová
veličina a proč mohou existovat poločíselné hodnoty, které orbitální moment
mít nemůže.
\end{vysvetleni}

\subsection{Základní členění částic}

Se znalostí spinu můžeme částice roztřídit dvěma nezávislými způsoby.

\textbf{Podle spinu} (určuje statistiku – Pauliho princip vs.\ Boseho–Einsteinova statistika):
\begin{itemize}
  \item \kw{fermiony} – poločíselný spin \(\left(\tfrac12, \tfrac32, \dots\right)\):
  \(p\), \(n\), \(e^-\), \(\nu_e\) (spin \(\tfrac12\)); \(\Delta^{++}\), \(\Omega^-\)
  (spin \(\tfrac32\));
  \item \kw{bosony} – celočíselný spin: \(\gamma\), \(\pi\), \dots
\end{itemize}
\vspace{-4pt}\noindent\rule{2.5cm}{0.4pt}

\textbf{Podle interakcí}:
\begin{itemize}
  \item \kw{hadrony} – interagují mj.\ silně: \(p\), \(n\), \(\pi\), \dots;
  dělí se dále na
  \begin{itemize}
    \item \kw{baryony} (fermiony) – např.\ \(p, n, \Delta, \Omega\),
    \item \kw{mezony} (bosony) – např.\ \(\pi\);
  \end{itemize}
  \item \kw{leptony} – interagují slabě, elektromagneticky (\emph{ne} silně):
  \(e^-\), \(\mu^-\), \(\nu_e\), \dots\ \(\to\) fermiony se spinem \(\tfrac12\).
\end{itemize}
Obě dělení se kříží: např.\ proton je hadron i fermion, pion hadron
i boson, elektron lepton i fermion. V kvarkovém modelu (později) se ukáže,
že baryony jsou složeny ze tří kvarků (proto poločíselný spin) a mezony
z kvarku a antikvarku (celočíselný spin).
